\GXNSFBodyText

\noindent\textbf{示例项目概述}\par
本模板以“平台化传播中公众人物跨媒介职业韧性的多模态时序建模”为例，展示广西自然科学基金面上项目正文的叙事、图表与公式组织。研究对象为佐佐木希（佐々木希 / Nozomi Sasaki）的公开职业资料；正式申请时，应以申请人的真实主题、数据与前期证据替换本示例。

佐佐木希的公开职业轨迹跨越杂志、广告、影视与自营视频等媒介，且阶段转换较为清晰。TopCoat 官方页面可核对其基本信息[1]；2024 年开设官方 YouTube 频道的节点亦有公开新闻稿可查[2]。这些可复核节点为观察媒介状态如何随时间变化提供了合适的演示材料。

\noindent\textbf{研究背景与科学意义}\par

\noindent\textbf{平台化传播形成复杂动态系统}\par
公众人物的职业发展不是作品数量或曝光频次的简单累加，而是在平台规则、内容形态、受众反馈和外部事件共同作用下不断调整的过程。纸媒、影视、综艺、广告、社交平台和自营视频频道拥有不同的信息扩散结构；主体进入新媒介后，既会延续既有形象资本，也会受到算法推荐、互动节奏与平台治理的重新塑形。融合文化研究指出内容会在多种媒介间流动[5]，自我呈现理论则提示公众形象会随情境持续调整[4]。因此，职业路径可被理解为多层媒介网络上的状态迁移过程。

\noindent\textbf{职业韧性需要可计算定义}\par
现有讨论常以转型成功、形象修复或重新出发概括职业变化，却很少区分短期热度反弹与结构性恢复。本示例将职业韧性界定为：主体在受到外部扰动后，仍能维持核心身份识别，并逐步恢复跨平台内容供给和受众连接的能力。恢复速度、媒介多样性、叙事一致性和波动稳定性共同构成可观测维度，据此比较不同平台组合下的恢复路径。

\noindent\textbf{广西区域场景的延伸价值}\par
广西面向东盟，是跨语言数字内容流动和区域文化传播的重要节点。本示例以日本公开人物资料贯通完整方法，并将多语言实体对齐、平台校准和扰动窗口构造等通用组件置于广西—东盟文化内容 IP 或区域品牌的固定观察窗中作探索性测试。由此，项目既关注多层媒介网络中的状态迁移与扰动恢复，也明确区分人物职业机制与区域文化传播场景的分析边界。

\noindent\textbf{国内外研究现状与不足}\par

\noindent\textbf{职业阶段研究缺少平台耦合机制}\par
经典职业发展研究将个体发展描述为探索、建立、维持和再适应等阶段[3]，为长期轨迹分析提供了基本框架。然而，访谈和横截面问卷难以同时呈现平台迁移、角色叠加与突发事件的交互影响。对于跨媒介职业样本，仅划分“模特期—影视期—平台期”仍难解释状态为何转换、何时转换，以及转换后能否保持稳定。

\noindent\textbf{计算社会科学提供方法但存在割裂}\par
计算社会科学强调利用大规模数字痕迹研究社会行为。文本情感、主题模型、图像表征和社会网络分析已能分别描述内容语义、视觉风格与传播结构，但多模态特征常被独立处理；不同平台的时间粒度、缺失机制和反馈口径也不一致。如果直接拼接特征，容易把平台差异误判为主体变化，并把相关性误写成因果关系。

\noindent\textbf{韧性研究缺少可反驳的边界条件}\par
韧性概念在生态、工程和组织研究中均强调系统受扰后的维持与恢复，但应用到公众人物职业系统时，常以少数高热度节点替代完整时间过程。现有研究较少回答三个问题：扰动发生前的基线如何确定；恢复是回到旧状态还是进入新稳定状态；不同平台间的替代与协同如何影响恢复。缺少这些边界条件，会使“韧性”退化为事后叙事。

\noindent\textbf{本项目的切入点}\par

\noindent\textbf{从公开事件链转向多层状态空间}\par
本示例拟将公开资料整理为“时间—平台—内容—事件—反馈”五元记录，以隐状态刻画职业阶段，以多层网络描述媒介连接，以外部事件界定扰动。研究比较多种分段和状态转移解释，并在替代指标、不同时间窗口与对照样本中检验结论是否稳定，从而避免把单一事件直接等同于职业变化的原因。

\noindent\textbf{从单案例演示转向可迁移方法}\par
佐佐木希案例用于贯通公开事实核验、事件编码、多模态表征、状态识别和韧性估计等环节；模型训练与稳健性验证则引入同期公开人物对照样本。研究最终形成可复用的方法框架及其适用边界，为数字平台中的职业轨迹研究提供可检验的分析路径。

\noindent\textbf{主要参考文献}\par

\noindent\textbf{公开资料与理论方法文献}\par
\begin{enumerate}[label={\textup{[\arabic*]}},leftmargin=*,itemsep=0pt]
  \item TopCoat. 佐佐木希官方资料[EB/OL]. \href{https://topcoat.co.jp/nozomi_sasaki}{TopCoat 艺人页面}.
  \item 株式会社ブルズ. 佐佐木希开设 YouTube 官方频道[EB/OL]. PR TIMES, 2024-07-05. \href{https://prtimes.jp/main/html/rd/p/000000020.000059044.html}{PR TIMES 新闻稿}.
  \item Super D E. The Psychology of Careers[M]. New York: Harper \& Brothers, 1957.
  \item Goffman E. The Presentation of Self in Everyday Life[M]. New York: Doubleday, 1959.
  \item Jenkins H. Convergence Culture: Where Old and New Media Collide[M]. New York: New York University Press, 2006.
  \item Lazer D, Pentland A, Adamic L, et al. Computational social science[J]. Science, 2009, 323(5915): 721--723.
  \item Hamilton W L, Leskovec J, Jurafsky D. Diachronic word embeddings reveal statistical laws of semantic change[C]. Proceedings of ACL, 2016: 1489--1501.
  \item Pearl J. Causality: Models, Reasoning, and Inference[M]. 2nd ed. Cambridge: Cambridge University Press, 2009.
\end{enumerate}
